\chapter{Introducción y Objetivos}

\section{Introducción}

La predicción de los resultados electorales ha sido históricamente uno de los desafíos más complejos dentro de las ciencias sociales y la estadística aplicada. En el contexto español, el sistema de representación proporcional mediante la Ley D'Hondt y la división en 52 circunscripciones provinciales añade una capa de complejidad técnica que requiere una precisión extrema en la estimación del voto a nivel local. Tradicionalmente, este análisis se ha basado en encuestas de intención de voto y sondeos a pie de urna (\textit{exit polls}). Sin embargo, estas metodologías enfrentan una crisis de fiabilidad creciente debido a fenómenos como el sesgo de respuesta, el aumento de los votantes indecisos y el denominado "voto oculto".

Este trabajo parte de la premisa de que el comportamiento electoral no es un fenómeno aislado de carácter puramente emocional, sino que responde de manera endógena a las condiciones materiales y socioeconómicas de los ciudadanos. La hipótesis central sostiene que los indicadores macroeconómicos y demográficos —tales como la renta per cápita, los niveles de desempleo, la desigualdad medida a través del Índice de Gini y la densidad de población— constituyen un rastro estadístico objetivo y medible. Al contrario que las encuestas, estos datos no dependen de la voluntad de comunicación del sujeto, sino de hechos económicos registrados oficialmente.

Mediante el uso de técnicas avanzadas de \textit{Data Engineering} y algoritmos de \textit{Machine Learning}, este proyecto propone un cambio de paradigma: la transición de una predicción basada en la opinión a una inferencia basada en la realidad socioeconómica. Se busca identificar patrones no lineales que vinculen la precariedad o la bonanza económica con tendencias ideológicas específicas, permitiendo modelizar la composición del Congreso de los Diputados sin recurrir a la consulta directa.

\section{Objetivos}

\subsection{Objetivo General}
Desarrollar y validar un sistema de inferencia predictiva que, mediante el análisis de indicadores socioeconómicos municipales y el uso de algoritmos de aprendizaje supervisado, sea capaz de estimar el reparto de los 350 escaños del Congreso de los Diputados en las elecciones generales españolas.

\subsection{Objetivos Específicos}
\begin{enumerate}
    \item \textbf{Ingesta y Fusión de Datos Heterogéneos:} Diseñar un pipeline de datos para la extracción y limpieza de fuentes oficiales (INE, SEPE, IGN) y su integración en un dataset único a nivel municipal que abarque el periodo 2016-2023.
    \item \textbf{Análisis Exploratorio y Correlacional:} Identificar las variables socioeconómicas con mayor poder predictivo y analizar su impacto en el sesgo de voto mediante técnicas de análisis multivariante.
    \item \textbf{Modelado de Inferencia Predictiva:} Implementar y optimizar modelos de \textit{Gradient Boosting} que capturen las relaciones complejas entre el entorno económico y la intención de voto, integrando componentes de influencia espacial.
    \item \textbf{Desarrollo del Motor de Simulación Electoral:} Programar un módulo lógico que traduzca las predicciones de voto municipal en escaños reales, aplicando con rigor la normativa electoral vigente (barreras del 3\% y Ley D'Hondt).
    \item \textbf{Validación y Backtesting:} Evaluar la precisión del modelo utilizando los resultados de las elecciones de 2023 como conjunto de prueba frente al entrenamiento basado en datos históricos anteriores.
\end{enumerate}
